% Fragmento público generado desde propietarios íntegros.
% Residencia: 52.2.3.
% Fuente: ../incoming/radion_monografia_20260827/manuscrito/sections/03_cierre_exacto.tex:104-253
\subsubsection{Cocycle de transition entre les sections directe et torsionnelle}

\paragraph{Définition opératoire de la retenue \(\varepsilon_R\)}

\begin{definition}[Retenue réelle du radion]
La sortie unitaire du radion est le cocycle de transition
\begin{equation}
\epsone:=S_T-S_E
=\Phi_T-D_E
=\frac{20}{81}\Delta_4-operatorname{rem}_1(S_E).
\label{eq:epsilon-def}
\end{equation}
\end{definition}

L'équation contient une soustraction, mais ce n'est pas un ajustement à la cible. Ses deux
opérandes \(S_T\) et \(S_E\) sont générés séparément. L'indépendance
démontrée est une indépendance par rapport à \(180/\pi\) et à un
\(\varepsilon\) cible ; on n'affirme pas d'indépendance algébrique entre une
différence et ses termes.

\begin{theorem}[Clôture exacte unitaire]
Avec \(S_E\), \(\Phi_T\) et \(\epsone\) définis par
\eqref{eq:seccion-electrica}, \eqref{eq:delta4}, \eqref{eq:rho-twoway} et
\eqref{eq:epsilon-def}, on a exactement
\begin{equation}
\boxed{1=S_E-\Phi_T+\epsone.}
\label{eq:cierre-unitario}
\end{equation}
\end{theorem}
\begin{proof}
Par \eqref{eq:SE-intervalo}, \(D_E=S_E-1\). Par définition,
\(\epsone=\Phi_T-D_E\). Alors
\[
S_E-\Phi_T+\epsone
=S_E-\Phi_T+\Phi_T-(S_E-1)=1.
\]
La force de la preuve ne réside pas dans cette annulation élémentaire, mais dans le fait que la
généalogie préalable fixe les deux opérandes sans utiliser la clôture comme entrée.
\end{proof}

\paragraph{Publication de la clôture dans la carte angulaire}

La carte du tour multiplie une fraction unitaire par
\(360/T_+=180/\pih\). En particulier,
\[
\epsR^\circ=\frac{360^\circ}{T_+}\epsone.
\]
Multiplier \eqref{eq:cierre-unitario} par \(360^\circ/T_+\) et utiliser
\[
\frac{360}{T_+}S_E
=360\left(\frac5{36}+\frac{25}{9}\alphah\right)
=50+1000\alphah=50+A
\]
produit la clôture angulaire.

\begin{theorem}[Équation exacte du radion]
\BigRadionEquation
\end{theorem}

L'équation dit quelque chose de plus précis que \enquote{le radian est
\(57.2957\ldots{}^\circ\)}. Elle décompose cette publication en quatre opérations
avec démonstration de référence :
\begin{center}
\begin{tabularx}{.96\textwidth}{>{\bfseries}l X}
\toprule
\(50^\circ\) & deuxième phase de la roue ; charnière critique dans une carte déclarée ;\\
\(A^\circ\) & coordonnée commune de la publication parangulaire ;\\
\(-\frac{20}{81}\Delta_4^\circ\) & transport orienté de la torsion globale ;\\
\(+\epsR^\circ\) & retenue réelle qui raccorde les deux relèvements.\\
\bottomrule
\end{tabularx}
\end{center}

\paragraph{Convergence selon la profondeur du raffinement}

À la profondeur \(N\), le générateur fournit des cylindres rationnels
\(p_N\to\pih\) et \(e_N\to e_{\HMT}\). On définit
\[
S_{E,N}=2p_N\left(\frac5{36}+\frac{25}{9}\alphah\right),
\]
\[
\Delta_{4,N}=\frac{p_N-e_N\log p_N}{270}
\left(1+\frac{p_N}{729}\right),\qquad
S_{T,N}=1+\frac{20}{81}\Delta_{4,N},
\]
et
\[
\mathcal E_N=S_{T,N}-S_{E,N}.
\]
La continuité de la fonctionnelle en \(3<p<4\), \(2<e<3\), jointe aux
dérivées monotones, transporte chaque couple de cylindres vers un intervalle
certifié. À la profondeur de \(45\) blocs, sa largeur est inférieure à
\(4.81\times10^{-130}\).

Chaque retour de neuf niveaux contient six trits par niveau. C'est pourquoi le taux
de raffinement est
\begin{equation}
729^{-9}=3^{-54}.
\label{eq:contraccion-54}
\end{equation}
Ce chiffre contrôle la largeur des intervalles ; ce n'est pas la valeur de
\(\epsR\).

La soustraction de triades stabilise
\begin{center}
\ttfamily
000|000|000|896|802|822|044|562|981|999|050|695|020|651|286|413
\end{center}
et le préfixe de retenues
\begin{center}
\ttfamily
0,0,0,0,0,-1,0,-1,-1,-1,0,-1,0.
\end{center}
Tel est le mécanisme concret de mémoire décimale : non une queue choisie, mais
la récurrence \eqref{eq:subacarreo} appliquée à des opérandes préalablement
construits.

\paragraph{Évaluation des valeurs canoniques}

La sortie coinductive est
\begin{align}
\epsone
&=8.9680282204456298199905069502065128641372736205009\ldots
\times10^{-10},\\
\epsR^\circ
&=5.1383016758575282071685164622645947489908574400054\ldots
\times10^{-8}\,{}^\circ.
\label{eq:epsilon-valores}
\end{align}
La publication complète donne
\[
\frac{180^\circ}{\pih}
=57.2957795130823208767981548141051703324054724665643\ldots{}^\circ.
\]

\begin{exactbox}
\begin{align*}
50^\circ+A^\circ
&=57.29735256928380099728510547238066299956\ldots{}^\circ,\\
\frac{20}{81}\Delta_4^\circ
&=0.00157310758449687906223272996065728980465\ldots{}^\circ,\\
50^\circ+A^\circ-\frac{20}{81}\Delta_4^\circ
&=57.29577946169930411822287274242000570976\ldots{}^\circ,\\
\epsR^\circ
&=0.00000005138301675857528207168516462264594749\ldots{}^\circ,\\
\text{somme finale}
&=57.29577951308232087679815481410517033241\ldots{}^\circ.
\end{align*}
\end{exactbox}

